\documentclass[11pt]{article}
\usepackage[margin=1in]{geometry}
\usepackage{setspace}
\usepackage{enumitem}
\usepackage{amsmath}
\usepackage{amssymb}
\usepackage{fancyhdr}
\IfFileExists{lastpage.sty}{\usepackage{lastpage}}{}
\usepackage{hyperref}

\setlength{\parindent}{0pt}
\setlength{\parskip}{0.6em}
\setlist[itemize]{leftmargin=*}
\setlist[enumerate]{leftmargin=*}

\newcommand{\answerlines}[1]{\vspace{#1\baselineskip}}

\pagestyle{fancy}
\fancyhf{}
\IfFileExists{lastpage.sty}{
  \fancyfoot[C]{Financial Data Science II\ \ \ Updated March 17, 2026\ \ \ Page \thepage\ of \pageref{LastPage}}
}{
  \fancyfoot[C]{Financial Data Science II\ \ \ Updated March 17, 2026\ \ \ Page \thepage}
}

\begin{document}

\begin{center}
{\Large \textbf{Part 23: Feature Selection and Temporal Validation}}
\end{center}

Part 22 considered the construction of candidate predictors. The next task is
to determine which of those predictors merit further attention.

This is not a mechanical step. A feature may appear useful in sample and fail
out of sample. A feature may also be statistically related to the target and
yet be infeasible in real time.

For that reason, feature selection is best viewed as a sequence of screens
rather than as a single algorithm.

\textbf{Exercise.} Explain the difference between the following statements:

\begin{enumerate}
  \item a feature is correlated with the target in sample
  \item a feature is useful in a predictive model
\end{enumerate}

\answerlines{5}

\newpage

\section*{Basic Terms}

A \textbf{screening rule} is a preliminary rule used to reduce a large
candidate set. A \textbf{selection rule} determines which features are retained
for modelling. \textbf{Redundancy} refers to overlap in the information carried
by different features. \textbf{Stability} refers to whether a feature behaves
similarly across periods or regimes. An \textbf{ablation study} compares a
baseline specification to an expanded one in order to assess the incremental
value of a feature or feature family.

Finally, we distinguish between a \textbf{ranking statistic} and a
\textbf{modelling objective}. For example, a feature may be ranked by a
univariate $F$-score or by mutual information, while the final predictive model
is judged by out-of-sample RMSE, $R^2$, AUC, or some other metric.

\textbf{Exercise.} Why is it possible for two features to rank highly on their
own but add little value when used together in the same model?

\answerlines{6}

\newpage

\section*{Why Feature Selection is Needed}

Feature selection matters for at least four reasons.

\begin{enumerate}
  \item Large feature sets can lead to unstable fitted models.
  \item Many features are redundant transformations of the same underlying signal.
  \item Financial relationships vary over time, so a feature that works in one period may fail in another.
  \item Some features are costly to maintain or impossible to compute in a production setting.
\end{enumerate}

The goal is not to maximize the number of features. The goal is to retain the
subset that is informative, stable, and usable.

\section*{A Selection Framework}

A practical sequence is the following.

\begin{enumerate}
  \item Remove features that are infeasible or contaminated by leakage.
  \item Remove features with too many missing values or too little variation.
  \item Rank the remaining features using simple univariate measures.
  \item Check whether the ranking is stable across time.
  \item Evaluate incremental predictive value using walk-forward validation.
\end{enumerate}

\section*{Univariate Screening}

In this Part we use two simple ranking statistics:

\begin{enumerate}
  \item \textbf{$F$-score:} emphasizes linear association with a continuous target.
  \item \textbf{Mutual information:} allows for more general dependence.
\end{enumerate}

A convenient summary is the \textbf{combined rank}, defined as the average of
the rank under each criterion. This is a screening device, not a final
objective function.

\newpage

\section*{Temporal Stability}

A feature that ranks well over the full sample may do so only because of one
favorable period. It is therefore useful to check whether a feature retains its
rank across rolling subsamples.

Stability is not a guarantee of predictive success, but unstable rankings are a
warning sign. This is particularly important in finance, where the underlying
relationship between predictors and outcomes is rarely fixed.

\section*{Ablation and Incremental Value}

A useful habit is to compare nested feature sets. One may begin with a simple
baseline, then add a new feature or feature family and ask whether the
out-of-sample performance improves.

This is preferable to moving immediately to a large feature set, because it
makes the contribution of the new information easier to interpret.

\newpage

\section*{What Should Be Reported?}

When presenting a selected feature set, it is useful to report at least the
following:

\begin{enumerate}
  \item the original candidate features
  \item the screening criteria
  \item the final selected subset
  \item the validation scheme
  \item the metric by which performance is judged
  \item the incremental contribution of the selected features relative to a simpler baseline
\end{enumerate}

This makes it easier to distinguish between a feature that is merely
interesting and a feature that is useful for a forecasting problem.

\section*{Exercises}

\textbf{Exercise 1.} Give an example of a feature that ranks well by mutual
information but might still be rejected for practical reasons.

\answerlines{6}

\textbf{Exercise 2.} Suppose that two features have nearly identical combined
rank values. What additional information would you want before choosing between
them?

\answerlines{6}

\textbf{Exercise 3.} In a LOB setting, why is it useful to compare a base
book-state model to a base-plus-OFI model rather than jumping immediately to
the full feature set?

\answerlines{6}

\textbf{Exercise 4.} Explain why a feature-selection rule based on a shuffled
train/test split is not appropriate in this setting.

\answerlines{6}

\section*{References}

\begin{enumerate}
  \item Cont, Kukanov, and Stoikov (2013), \textit{The Price Impact of Order Book Events}.
  \item Guyon and Elisseeff (2003), \textit{An Introduction to Variable and Feature Selection}.
  \item Lopez de Prado (2018), \textit{Advances in Financial Machine Learning}.
  \item LOBSTER documentation and sample files.
\end{enumerate}

\end{document}
